\section{Introduction}\label{sec:introduction}

How can one tell that two matrices generate a free group? By definition, one
has to check that no nonempty reduced word in them is the identity. There are
infinitely many words, so this cannot be done one word at a time.

The ping-pong lemma replaces the words by a picture. Let the group act on a set,
and find four disjoint pieces of that set, one for each of the letters $a$,
$a^{-1}$, $b$, $b^{-1}$. Suppose each letter throws the pieces of the other
letters (except its own inverse) into its own piece. Then a reduced word
bounces a point from piece to piece, like a ping-pong ball, and the point ends
up in the piece of the first letter of the word. If it started somewhere else,
it has moved, so the word is not the identity.

The argument is old. It is usually traced back to Klein~\cite{klein-1883}, who
combined two groups acting on the sphere into their free product when each
group moves the other's fundamental region inside its own; we follow de la
Harpe~\cite{delaharpe-2000} for this attribution. The model example is a
Schottky group, generated by M\"obius transformations that pair off disjoint
circles in the plane (Section~\ref{sec:schottky}). In 1947 Sanov~\cite{sanov-1947}
showed that the integer matrices
\[
  A = \mat{1}{2}{0}{1}, \qquad B = \mat{1}{0}{2}{1}
\]
generate a free group. The same holds for the family
\[
  A_t = \mat{1}{t}{0}{1}, \qquad B_t = \mat{1}{0}{t}{1}, \qquad \abs{t} \ge 2,
\]
which is usually credited to Brenner~\cite{brenner-1955}. In 1972 Tits~\cite{tits-1972} made ping-pong the
engine of a structure theorem: a finitely generated linear group either
contains a free subgroup of rank~$2$ or has a solvable subgroup of finite
index.

\subsection*{What this paper does}
Section~\ref{sec:lemma} proves the lemma in two forms: for free groups
(Theorem~\ref{thm:pingpong}, the main statement of this project) and for free
products (Theorem~\ref{thm:tabletennis}). Section~\ref{sec:sanov} proves
Sanov's theorem and its extension to $\abs{t} \ge 2$, and explains why the
argument must fail at $t = 1$, where the group is $\SL_2(\Z)$. Sections~\ref{sec:schottky}
and~\ref{sec:tits} treat Schottky groups and the Tits alternative,
Section~\ref{sec:open} describes an open question, and
Section~\ref{sec:project} points to the Lean files and computations that go
with this paper. For more, see de la Harpe's book~\cite{delaharpe-2000},
where the lemma is called the table-tennis lemma, and L\"oh's
textbook~\cite{loh-2017}, a friendly first read.
